\originalchapter{习题课六: 磁力与力矩}
题源为\href{https://ocw.mit.edu/courses/8-02-physics-ii-electricity-and-magnetism-spring-2007/resources/solving06/}{Problem Solving 6}, PDF第1--7页(原页脚沿用Solving 7). Q1--Q10及轨迹考试练习题全部保留. 原图的角$\theta$从竖直线量起, 不是从水平线量起. 所有解答为AI独立解答.

\begin{exercise}
质谱仪. 离子源从原子剥去一个电子, 正离子初速近零, 向上经过电势差加速, 进入匀强磁场. 它在上方偏转区弯向右侧计数器, 两孔水平间距$D$.
Q1加速器上方或下方哪端电势应更高?
Q2磁场应朝纸内还是纸外?
Q3进入偏转区时动能是多少?
Q4说明轨迹, 推导把质量$m$离子送进右孔所需$B$.
Q5若$B=1\,\mathrm T,D=20\,\mathrm{cm}$, 单电离碳12约需多少加速电压? 依原件近似每个核子静能约$1\,\mathrm{GeV}$.
\end{exercise}
\begin{solution}
Q1. 下方电势高, 正离子沿电场向上加速. 定义$V_a=V_{\rm bottom}-V_{\rm top}>0$.
Q2. 初速度向上, 初始弯向右, 所以$\uv y\times\bb$向右, $\bb$应指向纸外.
Q3. $K=qV_a$.
Q4. 磁力不改变速率, 半圆半径$D/2$; $v=\sqrt{2qV_a/m}$, 故$B=2mv/(qD)=\sqrt{8mV_a/q}/D$. 反解$V_a=qB^2D^2/(8m)$.
Q5. 使用$m\simeq12\,\mathrm{GeV}/c^2\simeq2.14\times10^{-26}\,\mathrm{kg}$、$q=e$, 得$V_a\simeq3.74\times10^4\,\mathrm V$. 这是按题设核子质量近似计算, 不等于使用真实碳原子质量的精密质谱值; 本册附带数值脚本实际复算该值.
\end{solution}

\begin{exercise}
线圈支承. 边长$\ell$方线圈绕水平上边转动, 匀强磁场竖直向下. 侧视图中从竖直向下方向到线圈边的夹角为$\theta$. 线圈电流$I$.
Q6求磁力矩大小.
Q7从上方看电流顺时针还是逆时针, 才能对抗重力使线圈抬起?
Q8线圈近乎无质量, 边长$1\,\mathrm m$, 一个$20\,\mathrm{kg}$孩子挂在下边, $I=100\,\mathrm A$, 要支承在$45^\circ$需多大$B$?
Q9小扰动会放大还是恢复? 若不稳定估时间, 若稳定求小振动频率. 可用三角展开, 孩子距轴$\ell$, 转动惯量$m\ell^2$.
Q10改为支承在$60^\circ$, 竖直场与选择适当方向的水平场哪种所需$B$更小?
\end{exercise}
\begin{solution}
Q6. 线圈法向与竖直方向的夹角为$90^\circ-\theta$或补角, 故$|\tau_B|=I\ell^2B\cos\theta$.
Q7. 从上方看顺时针使磁矩具有向下分量, 与向下场对齐会降低势能, 因而使线圈向较大的$\theta$抬起.
Q8. 增大$\theta$为正, 重力矩$-mg\ell\sin\theta$, 平衡为$I\ell^2B\cos\theta=mg\ell\sin\theta$. 所以$B=[mg/(I\ell)]\tan\theta$, 在$45^\circ$约$1.96\,\mathrm T$.
Q9. 总力矩$\tau=I\ell^2B\cos\theta-mg\ell\sin\theta$. 写$\theta=\pi/4+\delta$, 用平衡条件$I\ell^2B=mg\ell$得$\tau\simeq-\sqrt2mg\ell\delta$. 因而$m\ell^2\ddot\delta=-\sqrt2mg\ell\delta$, 是稳定恢复振动, $\omega=\sqrt{\sqrt2g/\ell}$, 频率$f=\omega/(2\pi)\simeq0.593\,\mathrm{Hz}$. 若线圈本身质量不忽略, 要把其重力矩与惯量加进来.
Q10. 竖直场需$B_v=\sqrt3mg/(I\ell)$; 选择水平场让力矩抬起时, 磁力矩为$I\ell^2B_h\sin\theta$, 所需$B_h=mg/(I\ell)$, 比较可知水平场更小. 但该理想模型中$B_h$恰好取平衡值时两力矩对任何角都抵消, 没有竖直场的恢复刚度. “所需场小”不等于稳定性更好.
\begin{center}\begin{tikzpicture}[>=Stealth]
\fill (0,0)circle(2pt)node[above]{转轴};\draw[dashed](0,0)--(0,-2.6);
\draw[thick](0,0)--(1.6,-1.6)node[midway,right]{$\ell$};
\draw[->](0,-.7)arc(-90:-45:.7)node[midway,below]{$\theta$};
\draw[->](1.6,-1.6)--(1.6,-2.5)node[right]{$mg$};
\draw[->](-1,.2)--(-1,-1.2)node[left]{$\bb$};
\end{tikzpicture}\end{center}
\end{solution}

\begin{exercise}
考试练习题. $q>0$、质量$m$的粒子在原点以速率$V$沿$+y$出发. 原图轨迹先在$y>0$沿半径$R$的上半圆弯向右, 再在$y<0$沿半径$R/2$的下半圆回到$x$轴上的$P$, 速率始终不变.
(a)上半平面磁场朝纸内还是外?
(b)求其大小.
(c)下半平面的方向和大小.
(d)从原点到$P$所需时间.
\end{exercise}
\begin{solution}
(a)向纸外, 初始磁力向右. (b)$B_+=mV/(qR)$.
(c)刚进入下半平面时速度向下而曲率朝右, 所以下部磁场向纸内, 大小$B_-=2mV/(qR)$.
(d)上段弧长$\pi R$, 下段弧长$\pi R/2$, 总时间$3\pi R/(2V)$. 原图$P$位于$x=3R$, 不是上半圆首次落回$x$轴的$x=2R$点.
\end{solution}
